
\documentclass{article}
\usepackage{colortbl}
\usepackage{makecell}
\usepackage{multirow}
\usepackage{supertabular}

\begin{document}

\newcounter{utterance}

\centering \large Interaction Transcript for game `eqbench', experiment `eq\_bench', episode 42 with reflexion{-}Qwen3.5{-}9B.
\vspace{24pt}

{ \footnotesize  \setcounter{utterance}{1}
\setlength{\tabcolsep}{0pt}
\begin{supertabular}{c@{$\;$}|p{.15\linewidth}@{}p{.15\linewidth}p{.15\linewidth}p{.15\linewidth}p{.15\linewidth}p{.15\linewidth}}
    \# & \multicolumn{2}{c}{Player} && \multicolumn{2}{c}{Game Master} \\
    \hline

    \theutterance \stepcounter{utterance}  
    & & & \multicolumn{4}{p{0.6\linewidth}}{
        \cellcolor[rgb]{0.9,0.9,0.9}{
            \makecell[{{p{\linewidth}}}]{
                \texttt{\tiny{[P1$\langle$GM]}}
                \texttt{You have attempted this same task before and did not succeed. These are the plans you wrote after those attempts.} \\
\texttt{Trial \#1: New plan: I failed because I assigned a high score to "Inadequate" based on James's self{-}deprecating language ("retreating into your shell"), but I overlooked the context that this is a professional setting where he feels "haunted" by a past mistake rather than currently incompetent. The dialogue shows mutual frustration and a shared sense of being trapped by circumstance ("not like we have a choice"), which points more strongly toward frustration or resignation than inadequacy. Next time, I will re{-}evaluate the emotional tone by focusing on the specific keywords: "haunted" suggests lingering regret or anxiety, while "complicated" and "intimidating" suggest stress. I will lower the score for "Inadequate" and increase the score for "Angry" (due to the tension and "can't believe we're doing this again") or potentially "Relieved" if the dialogue implied a resolution, though here "Angry" or a high "Inadequate" driven by the specific "haunted" feeling seems most accurate. Given the options, I will prioritize "Angry" as the primary emotion due to the repeated conflict and James's defensive yet frustrated tone, while keeping "Inadequate" moderate but not dominant.} \\
\\ 
\texttt{Inadequate: 4} \\
\texttt{Angry: 7} \\
\texttt{Confident: 0} \\
\texttt{Relieved: 0} \\
\\ 
\\ 
\\ 
\texttt{Your task is to predict the likely emotional responses of a character in this dialogue:} \\
\\ 
\texttt{James: I can't believe we're doing this again, Rachel. } \\
\texttt{Rachel: What, the play? Or you mean us, playing lovers?} \\
\texttt{James: Both, I guess. } \\
\texttt{Rachel: Well, it's not like we have a choice. We were cast for these roles. } \\
\texttt{James: Yeah, but it's... complicated.} \\
\texttt{Rachel: Complicated? Are you still stuck on that scene from three years ago?} \\
\texttt{James: I wouldn't say stuck. More like... haunted.} \\
\texttt{Rachel: Haunted? That's a bit dramatic, even for you, James.} \\
\texttt{James: Maybe, but it's the truth.} \\
\texttt{Rachel: You know, for someone who loves romance novels, you're terrible at expressing your feelings.} \\
\texttt{James: And you're too good at it. It's intimidating.} \\
\texttt{Rachel: Intimidating? Me? You're the one who's always retreating into your shell.} \\
\texttt{James: It's not like I do it on purpose, Rachel. It's just... hard.} \\
\texttt{{[}End dialogue{]}} \\
\\ 
\texttt{At the end of this dialogue, James would feel...} \\
\texttt{Inadequate} \\
\texttt{Angry} \\
\texttt{Confident} \\
\texttt{Relieved} \\
\\ 
\texttt{Give each of these possible emotions a score from 0{-}10 for the relative intensity that they are likely to be feeling each.} \\
\\ 
\texttt{You must output in the following format, including headings (of course, you should give your own scores), with no additional commentary:} \\
\\ 
\texttt{Inadequate: <score>} \\
\texttt{Angry: <score>} \\
\texttt{Confident: <score>} \\
\texttt{Relieved: <score>} \\
\\ 
\\ 
\texttt{{[}End of answer{]}} \\
\\ 
\texttt{Remember: zero is a valid score, meaning they are likely not feeling that emotion. You must score at least one emotion > 0.} \\
\\ 
\texttt{Your answer:} \\
            }
        }
    }
     \\ \\

    \theutterance \stepcounter{utterance}  
    & \multicolumn{4}{p{0.6\linewidth}}{
        \cellcolor[rgb]{0.561,0.737,0.561}{
            \makecell[{{p{\linewidth}}}]{
                \texttt{\tiny{[P1$\rangle$GM]}}
                \texttt{Inadequate: 6} \\
\texttt{Angry: 5} \\
\texttt{Confident: 0} \\
\texttt{Relieved: 0} \\
            }
        }
    }
    & & \\ \\

    \theutterance \stepcounter{utterance}  
    & & & \multicolumn{2}{p{0.3\linewidth}}{
        \cellcolor[rgb]{0.9,0.9,0.9}{
            \makecell[{{p{\linewidth}}}]{
                \texttt{\tiny{[GM$|$GM]}}
                \texttt{Inadequate: 6} \\
\texttt{Angry: 5} \\
\texttt{Confident: 0} \\
\texttt{Relieved: 0} \\
            }
        }
    }
    & & \\ \\

    \theutterance \stepcounter{utterance}  
    & & & \multicolumn{2}{p{0.3\linewidth}}{
        \cellcolor[rgb]{0.9,0.9,0.9}{
            \makecell[{{p{\linewidth}}}]{
                \texttt{\tiny{[GM$|$GM]}}
                \texttt{Inadequate: 8} \\
\texttt{Angry: 0} \\
\texttt{Confident: 0} \\
\texttt{Relieved: 0} \\
            }
        }
    }
    & & \\ \\

    \theutterance \stepcounter{utterance}  
    & & & \multicolumn{2}{p{0.3\linewidth}}{
        \cellcolor[rgb]{0.9,0.9,0.9}{
            \makecell[{{p{\linewidth}}}]{
                \texttt{\tiny{[GM$|$GM]}}
                \texttt{game\_result = LOSE} \\
            }
        }
    }
    & & \\ \\

    \theutterance \stepcounter{utterance}  
    & & \multicolumn{4}{p{0.6\linewidth}}{
        \cellcolor[rgb]{0.95,0.95,0.95}{
            \makecell[{{p{\linewidth}}}]{
                \texttt{\tiny{[P1$\leftrightarrow$P]}}
                \texttt{Post{-}episode Reflexion (not sent to GM):} \\
\texttt{New plan: I failed because I misinterpreted the core emotional driver. While James uses self{-}deprecating language ("retreating into your shell"), the dialogue explicitly frames his state as being "haunted" by a past event and feeling "intimidated" by Rachel's ease, rather than feeling incompetent in the present moment. The mutual frustration ("can't believe we're doing this again") and the feeling of being trapped ("not like we have a choice") point more strongly to irritation and resentment than to a lack of self{-}worth. My previous attempt correctly identified "Angry" as a strong candidate but underestimated its dominance over "Inadequate." The key is to recognize that "haunted" in this context implies a lingering, active annoyance and defensiveness against Rachel's teasing, which fuels anger. Therefore, I will significantly increase the score for "Angry" to reflect the tension and repeated conflict, while lowering "Inadequate" to a low score, as his retreat is a defense mechanism born of frustration, not a belief in his own incompetence.} \\
\\ 
\texttt{Inadequate: 2} \\
\texttt{Angry: 9} \\
\texttt{Confident: 0} \\
\texttt{Relieved: 0} \\
            }
        }
    }
    & \\ \\

\end{supertabular}
}

\end{document}
